\uuid{CrYL}
\titre{Changement de variables/éléments simples}
\auteur{hanani ; mzari}
\organisation{exo7}
\licence{CC BY-NC-SA 4.0}
\datecreate{2020-01-20}
\niveau{L1}
\module{Analyse}
\chapitre{Calcul d'intégrales}
\sousChapitre{Primitives diverses}
\difficulte{3}
\theme{Primitives, changement de variables, éléments simples}
\exo7qcmid{lille-2.primitives.047}

\contenu{
\begin{question}
\qclassification{127.11}
\qtags{motcle=changement de variables/éléments simples}
On se place dans l'intervalle $]-\pi,\pi[$ et on rappelle que $\displaystyle \cos x=\frac{1-\tan ^2(x/2)}{1+\tan ^2(x/2)}$ et que $\displaystyle \sin x=\frac{2\tan (x/2)}{1+\tan ^2(x/2)}$. Le changement de variable $t=\tan (x/2)$ donne :
\begin{answers}
\good{$\displaystyle \int \frac{\mathrm{d}x}{2+\cos x}=\int \frac{2\mathrm{d}t}{3+t^2}$.}
\good{$\displaystyle \int \frac{\mathrm{d}x}{2+\sin x}=\int \frac{\mathrm{d}t}{t^2+t+1}$.}
\bad{$\displaystyle \int \frac{\mathrm{d}x}{2+\cos x}=\frac{1}{\sqrt{3}}\arctan\left(\frac{t}{\sqrt{3}}\right)+k$, $k\in \Rr$.}
\bad{$\displaystyle \int \frac{\mathrm{d}x}{2+\sin x}=\frac{1}{\sqrt{2t+1}}\ln (t^2+t+1)+k$, $k\in \Rr$.}
\end{answers}
\vskip2mm
\begin{explanations}
Avec $t=\tan (x/2)$, on a : $\displaystyle \mathrm{d}x=\frac{2\mathrm{d}t}{1+t^2}$,
$$\frac{1}{2+\cos x}=\frac{1+t^2}{3+t^2}\quad \mbox{et}\quad \frac{1}{2+\sin x}=\frac{1+t^2}{2(t^2+t+1)}.$$
D'où
$$\int \frac{\mathrm{d}x}{2+\cos x}=\int \frac{2\mathrm{d}t}{3+t^2}=\frac{2}{\sqrt{3}}\arctan\left(\frac{t}{\sqrt{3}}\right)+k,\; k\in \Rr.$$
et
$$\int \frac{\mathrm{d}x}{2+\sin x}=\int \frac{\mathrm{d}t}{t^2+t+1}=\frac{2}{\sqrt{3}}\arctan\left(\frac{2t+1}{\sqrt{3}}\right)+k,\; k\in \Rr.$$
\end{explanations}
\end{question}
}
